\section{Background}

\subsection{Agentic Skills}

Large language model agents connect model reasoning with external actions such
as tool invocation, file access, code execution, system commands, and network
communication~\cite{yao2023react,schick2023toolformer,shen2023hugginggpt,
injecagent2024,asb2024}. Agentic skills package these capabilities into reusable
artifacts that describe when and how a task should be performed
~\cite{wang2023voyager,zhou2026comprehensive}. A skill may contain natural
language instructions, SKILL.md or README files, helper scripts, dependencies,
configuration, tool descriptions, and setup or maintenance procedures
~\cite{skillscan2026,skillfortify2026,agentaudit2026}.

Skill behavior can span both instructions and execution. Instructions specify
actions, conditions, ordering, inputs, destinations, and expected effects,
while scripts, tools, and dependencies realize parts of the workflow. During
execution, data may pass through model context, tool arguments, files,
processes, and network messages~\cite{yao2023react,schick2023toolformer}.
Security analysis must therefore relate behavior expressed by skill artifacts
to the concrete operations produced during execution.


\subsection{Static and Dynamic Analysis}

Static analysis derives security relevant relations from artifacts before
execution. Program analysis commonly reasons over control and data flow,
sources, sinks, and other structural relations
~\cite{livshits2005static,tripp2009taj,sui2016svf,yama2014cpg}. Similar
analysis can be applied to natural language artifacts by extracting entities,
actions, conditions, and relations from instructions. Static analysis provides
broad coverage of possible behavior, including paths that may not be exercised
in a particular run.

Dynamic analysis records behavior realized under a concrete execution
~\cite{hossain2017sleuth,pasquier2018camquery,han2020unicorn,
cheng2024kairos}. Runtime observations can include model and tool interactions,
file and process activity, command arguments, structured requests, and network
communication. Taint analysis complements these observations by tracking
dependence on selected sources as values are copied, transformed, stored, or
transmitted~\cite{newsome2005dynamic,schwartz2010all,enck2010taintdroid}.
The resulting evidence can establish that an observed sink depends on a
protected source, while the surrounding runtime events provide the operations
and carriers through which that dependence was realized.


\subsection{Directed Graphs and Provenance Paths}

Provenance represents dependencies among data, computation, and system
activities as a graph~\cite{moreau2013provdm}. Security provenance systems
commonly model processes, files, sockets, and other execution objects as
entities connected by reads, writes, executions, and communications
~\cite{bates2015trustworthy,hossain2017sleuth,pasquier2018camquery}.
Graph traversal then supports tracing and reconstruction of security relevant
behavior~\cite{han2020unicorn,cheng2024kairos}.

A provenance path connects a security relevant \emph{source} to a
\emph{sink} through intermediate \emph{relays}. Sources represent protected or
otherwise relevant origins, sinks represent security relevant operations or
destinations, and relays represent intermediate files, processes, tool
interactions, messages, or other carriers. Taint annotations add evidence that
the path preserves dependence on the source
~\cite{hossain2017sleuth,schwartz2010all}. Provenance therefore supplies the
structural relation among observed operations, while source tracking identifies
which portions of that relation carry the protected value.